\section{A local spinor obstruction to two horizon spheres}
\label{q3:local-spinor:section}

We give an independent proof of the two-sphere strictness conclusion.
It uses the complete static base from Section~\ref{q3:static:section}
and its end estimates, without the global classification in
Section~\ref{q3:global:section}. The reason is local: a conformal change makes
the base flat near each horizon, and the common surface gravity then
fixes the area of each sphere. Under equality, that area is already the
entire permitted horizon area. We obtain the local flatness by solving
Dirac equations on a family of complete metrics and taking a nonzero
parallel-spinor limit.

\begin{theorem}[Strictness for two spherical components]
\label{q3:local-spinor:strictness}
Let $(M,g,K)$ be smooth connected orientable three-dimensional initial
data without boundary, with $g$ complete and with finitely many
asymptotically flat ends outside a compact set. Assume on every end
\[
 g-\delta=O_2(r^{-q}),\qquad K=O_1(r^{-1-q}),\qquad q>\tfrac12,
\]
integrable constraint densities $\mu,|J|_g$, finite ADM charges, and
$\mu\geq |J|_g$ everywhere. Here
$16\pi\mu=R_g+(\operatorname{tr}_gK)^2-|K|_g^2$ and
$8\pi J=\operatorname{div}_g(K-(\operatorname{tr}_gK)g)$.
Let $S_1,S_2$ be disjoint smooth embedded two-sided spheres, and suppose
that the closure $\Omega$ of the connected exterior toward a designated
end has boundary $S=S_1\sqcup S_2$ and exactly that one end.

Orient $S$ into $\Omega$. Assume $\theta_+=H+\operatorname{tr}_S K=0$
on $S$, and assume every full enclosing cut has area at least
$A=|S_1|_g+|S_2|_g$. A full cut is the entire compact smooth frontier
of a connected end-containing outer domain, with all components and
coincident portions of $S$ counted and bounded complementary pockets
filled. Finally assume there is no such cut $\Gamma\ne S$ whose
components either equal an $S_i$ or lie wholly in the open exterior,
with at least one component in the open exterior and
$\theta_+(\Gamma)\leq0$ toward the designated end.
If that end has $E>|P|_\delta$, then
\[
 \sqrt{E^2-|P|_\delta^2}
    >\sqrt{\frac{|S_1|_g+|S_2|_g}{16\pi}}.
\]
\end{theorem}

The partial coincidence permitted in the last hypothesis excludes
replacing only one horizon component by a weakly trapped exterior
component while leaving the other in place. It is exactly the
outermostness used to localize the variational area sources.
The strict inequality is qualitative; it gives no uniform positive gap.

Write $m=\sqrt{E^2-|P|_\delta^2}>0$. The original exterior is complete
with its smooth boundary included: an intrinsic Cauchy sequence has a
limit in the closed subset $\Omega$ of the complete ambient manifold,
and a smooth half-ball chart at a boundary limit gives intrinsic
convergence. All constraint and end hypotheses restrict to $\Omega$.
Proposition~\ref{q3:intro:exterior-inequality} therefore applies.
Since $S$ itself is a full cut, outer area minimization identifies its
cut infimum with $A$, and the numerical inequality gives
$m\geq\sqrt{A/(16\pi)}$. It remains to exclude
\begin{equation}\label{q3:local-spinor:equality}
 A=16\pi m^2.
\end{equation}
Assume this equality for the remainder of the argument. We use the
original-data variational and static-base results only up to the
completion and harmonic-coordinate lemmas. The global classification
and hypersurface reconstruction play no role in the proof below.

\subsection{The complete metrics used in the limiting argument}

Proposition~\ref{q3:variation:adjoint} supplies the causal stationary
field on the original exterior with the same constant
\(\kappa=1/(4m)\) at both boundary components. The localization of the area sources in
Lemma~\ref{q3:variation:area-support} uses precisely the assumed
partial-coincidence outermostness. Lemma~\ref{q3:static:base-completion}
then supplies vanishing twist, a connected positive set
\(\mathcal P=\{N>0\}\), the complete static base and its two-copy
double. Its construction does not use the exclusion of the internal
null set or the Euclidean classification. In particular,
an internal null set may still give additional complete ends here.

Denote the double by \(W\), its metric by \(h_{\mathrm d}\), and
its signed lapse by \(\lambda_{\mathrm d}\). It is a connected
orientable smooth manifold without boundary. The metric and signed
lapse are \(C^2\), the metric is complete, and
\begin{equation}\label{q3:local-spinor:double-input}
 R_{h_{\mathrm d}}=0,\qquad
 \Delta_{h_{\mathrm d}}\lambda_{\mathrm d}=0,\qquad
 |\lambda_{\mathrm d}|<1.
\end{equation}
On the plus interior they equal the smooth fields \(h,\lambda\).
At each attached sphere, for the normal \(\nu_h\) into the plus side,
\begin{equation}\label{q3:local-spinor:horizon-input}
 \lambda=0,\qquad \II_h=0,\qquad
 h|_{TS_i}=g|_{TS_i},\qquad
 \nu_h\lambda=\kappa=\frac1{4m},\qquad i=1,2.
\end{equation}
To obtain complete scalar-flat metrics whose end energy tends to zero,
define, for \(0<t<1\),
\begin{equation}\label{q3:local-spinor:family}
 q_t=\left(\frac{1+t\lambda_{\mathrm d}}{1+t}\right)^4h_{\mathrm d},
 \qquad
 q_* =\left(\frac{1+\lambda_{\mathrm d}}2\right)^4h_{\mathrm d}.
\end{equation}
For fixed \(t<1\), its conformal factor satisfies
\[
 0<\frac{1-t}{1+t}
 <\frac{1+t\lambda_{\mathrm d}}{1+t}<1.
\]
Consequently \(q_t\) is complete. The three-dimensional conformal
scalar-curvature identity
\[
 R_{\phi^4h}=\phi^{-5}(-8\Delta_h\phi+R_h\phi)
\]
and Equation~\eqref{q3:local-spinor:double-input} make it scalar flat.
The identity holds across the joins by their matching second jets.
The metrics converge locally in \(C^2\) to the positive metric
\(q_*\). The lower comparison constant above depends on \(t\);
no completeness of \(q_*\) is used in this argument.

\subsection{Cancellation at the distinguished end}

Choose \(1/2<q_0<\min\{q,1\}\). The harmonic-coordinate proof of
Lemma~\ref{q3:static:asymptotics}, applied before any compactification
or global classification, gives on the plus end
\begin{equation}\label{q3:local-spinor:improved-end}
 \gamma=\lambda^2h=\delta+O_1(r^{-2q_0}),\qquad
 V=\log\lambda=\frac{a_\infty}{r}+O_1(r^{-2q_0})
\end{equation}
for a real constant \(a_\infty\). Indeed that proof uses
\(\epsilon=2q_0-1\), and its lapse expansion is stronger than the
one displayed here. Its coordinate construction starts with the original
two metric derivatives and then uses the static elliptic system to
bootstrap; no higher original falloff is added.

The spinor construction below needs an estimate uniform in the conformal
parameter. Fix \(t_0\in(0,1)\). On this end write
\[
 q_t=F_t(V)^4\gamma,\qquad
 F_t(v)=\frac{e^{-v/2}+t e^{v/2}}{1+t}.
\]
For \(t_0\le t\le1\) and small \(|v|\),
\[
 F_t(0)=1,\qquad F_t'(0)=\frac{t-1}{2(1+t)},\qquad
 |F_t'(v)|\le C((1-t)+|v|).
\]
Equation~\eqref{q3:local-spinor:improved-end} gives
\(|V|\le Cr^{-1}\), \(|\partial V|\le Cr^{-2}\), and
\(|\partial\gamma|\le Cr^{-1-2q_0}\). Hence, with \(q_1=q_*\),
\begin{equation}\label{q3:local-spinor:uniform-derivative}
 |\partial q_t|\le C\bigl((1-t)r^{-2}+r^{-1-2q_0}\bigr),
 \qquad t_0\le t\le1,
\end{equation}
where \(C\) is independent of \(t\). Here \(r^{-3}\le
r^{-1-2q_0}\) for \(r\ge1\). These metrics are uniformly
comparable to the Euclidean metric on a single fixed far end. More
precisely, the same calculation gives the cancellation
\begin{equation}\label{q3:local-spinor:leading-cancellation}
 (q_t)_{ij}=
 \left(1+\frac{2a_\infty(t-1)}{(1+t)r}\right)\delta_{ij}
       +O_1(r^{-2q_0}),
\end{equation}
uniformly in \(t\). The derivative estimate will control all end
errors in the spinor argument.

\subsection{Constructing spinors on the complete metrics}

We use the Dirac method of Witten~\cite{Witten1981}, constructing the
spinors on the complete metrics $q_t$ before taking the local limit.

\begin{lemma}[Spin structure and the compact energy identity]
\label{q3:local-spinor:structure}
The manifold $W$ admits a fixed spin structure.  For any metric
$q_t$ in Equation~\eqref{q3:local-spinor:family}, its spin connection
$\nabla^S_t$, Dirac operator $D_t$, and Clifford multiplication
$\mathfrak c_t$ satisfy
\begin{equation}\label{q3:local-spinor:lichnerowicz}
 \|D_t\zeta\|_{L^2(q_t)}^2
       =\|\nabla^S_t\zeta\|_{L^2(q_t)}^2
\end{equation}
for every compactly supported $H^1$ spinor $\zeta$.
\end{lemma}

\begin{proof}
Every singular two-cycle with coefficients in $\mathbb Z/2$ has
compact support and lies in the interior of a compact smooth
three-dimensional submanifold $C\subset W$ with boundary.
The closed orientable double of $C$ is parallelizable by
\cite[Theorem 1.1]{BenedettiLisca2018}.  Naturality therefore shows
that $w_2(TW)$ evaluates trivially on this cycle.  Duality between
cohomology and homology over the field $\mathbb Z/2$ gives
$w_2(TW)=0$.  Together with orientability this is the spin-structure
criterion; see \cite[Chapter II, Sections 1--2]{LawsonMichelsohn1989}.
This argument applies directly to the possibly noncompact $W$.

Fix one such spin structure.  Positive square roots of the changes
of metric identify the oriented orthonormal-frame bundles and their
spinor bundles.  We choose the lifts continuously in $t$; on compact
sets these identifications and the metric coefficients converge in
$C^2$ as $t\uparrow1$.  On the distinguished end the orthonormalized
coordinate frame has a spin lift because that end is simply
connected.  We choose these lifts consistently throughout the family.

We use the Clifford convention
$\mathfrak c_t(Y)\mathfrak c_t(Z)+\mathfrak c_t(Z)\mathfrak c_t(Y)
=-2q_t(Y,Z)$; Clifford multiplication by a real tangent vector is
skew-adjoint and has norm $|Y|_{q_t}$.  The local spin-curvature
formula and the curvature symmetries yield
\[
 D_t^2=(\nabla^S_t)^*\nabla^S_t+\tfrac14R_{q_t};
\]
the standard computation and formal adjoint identities are given in
\cite[Chapter~II, Sections~4--5 and~8]{LawsonMichelsohn1989}.
Integrating this identity for a compact test spinor proves
Equation~\eqref{q3:local-spinor:lichnerowicz} for smooth scalar-flat metrics.
For the present $C^2$ metric, smooth it in $C^2$ on a neighborhood
of the test support and use the positive-square-root bundle
identifications.  The connections and Dirac coefficients converge,
and the scalar curvatures converge uniformly to zero.  Passing to
the limit proves the identity for smooth compact test spinors.
Density in local $H^1$ then proves the stated form.  The joined
spheres are interior surfaces of the $C^2$ metric on $W$, so this
calculation has no boundary at a joined sphere.
\end{proof}

\begin{lemma}[Coercivity from one end]
\label{q3:local-spinor:coercivity}
Fix $t_0\in(0,1)$ and a sufficiently large end radius $r_0$.
For every compactly supported $H^1$ spinor $\zeta$ and
$t_0\le t<1$,
\begin{equation}\label{q3:local-spinor:hardy}
 \int_{r\ge r_0}r^{-2}|\zeta|^2\,dV_{q_t}
       \le C\|\nabla^S_t\zeta\|_{L^2(q_t)}^2.
\end{equation}
For every fixed compact $K\subset W$ there is also a constant
$C_K$, independent of $t$, such that
\begin{equation}\label{q3:local-spinor:local-coercivity}
 \|\zeta\|_{L^2(K,q_t)}
       \le C_K\|\nabla^S_t\zeta\|_{L^2(W,q_t)}.
\end{equation}
In fixed local bundle identifications the corresponding local
$H^1$ norm is bounded by the same global derivative norm.
\end{lemma}

\begin{proof}
For a compactly supported absolutely continuous function on a ray,
integration by parts gives
\[
 \int_{r_0}^{\infty}f^2\,dr
   =-r_0 f(r_0)^2-2\int_{r_0}^{\infty}rff'\,dr
   \le2\left(\int f^2\,dr\right)^{1/2}
          \left(\int r^2|f'|^2\,dr\right)^{1/2}.
\]
Consequently $\int f^2\le4\int r^2|f'|^2$; the inner boundary
term has the favorable sign without a prescribed inner trace.
Apply this to $f=|\zeta|$, integrate over directions, and use
uniform metric comparability from Equation~\eqref{q3:local-spinor:uniform-derivative}
and the Kato inequality
$|d|\zeta||_{q_t}\le|\nabla^S_t\zeta|_{q_t}$.
This proves Equation~\eqref{q3:local-spinor:hardy}, first for smooth
spinors and then by $H^1$ approximation.

Choose a fixed end annulus $\mathcal A$ and a relatively compact
connected smooth domain $G\subset W$ containing both $K$ and
$\mathcal A$.  The scalar Poincar\'e inequality with an anchored
mean gives
\[
 \|f\|_{L^2(G)}\le C_G
       \bigl(\|df\|_{L^2(G)}+\|f\|_{L^2(\mathcal A)}\bigr).
\]
For completeness, subtract the mean of $f$ on $G$ and apply the
usual Poincar\'e inequality; its restriction to $\mathcal A$
bounds that mean by the two terms on the right.
Take $f=|\zeta|$.  Equation~\eqref{q3:local-spinor:hardy} bounds its
annular term, and Kato bounds its derivative.  All norms on this
fixed domain are uniformly equivalent as $t\uparrow1$, since
$q_t\to q_*$ positively in $C^2$ there.  This proves
Equation~\eqref{q3:local-spinor:local-coercivity}.  In a fixed local
trivialization $\partial\zeta$ differs from
$\nabla^S_t\zeta$ by a uniformly bounded connection matrix times
$\zeta$, which gives the last assertion.
\end{proof}

For corrections with uniformly bounded derivative norm, the preceding
estimate gives local compactness and a uniform weighted bound on the
distinguished end. That bound will later prevent the corrections from
canceling a nonzero constant reference field on the whole end.

We take real parts of Hermitian inner products in the integral
pairings that follow.  For each fixed $t<1$, complete the compact
test spinors in the norm $\|\nabla^S_t\cdot\|_2$ and call the
resulting Hilbert space $\mathcal H_t$.
Lemma~\ref{q3:local-spinor:coercivity} identifies each element with a
well-defined local $H^1$ spinor, and extends both coercivity
inequalities to $\mathcal H_t$.  Equation~\eqref{q3:local-spinor:lichnerowicz}
then makes
\begin{equation}\label{q3:local-spinor:dirac-isometry}
 D_t:\mathcal H_t\longrightarrow L^2(W,q_t)
\end{equation}
an isometry with closed image.  The derivative and Dirac operator
of a completed element agree with its distributional derivatives,
by local $H^1$ convergence.

Choose a unit constant vector $\sigma\in\mathbb C^2$ in the coordinate
spin frame on the distinguished end and a smooth scalar cutoff
which equals one for $r\ge2r_0$ and vanishes before $r=r_0$.
Let $\psi_{0,t}$ be the resulting spinor on $W$, extended by zero
off this end neighborhood.  Its covariant derivative on the far
end is its spin connection acting on the fixed vector $\sigma$.
Thus Equation~\eqref{q3:local-spinor:uniform-derivative} implies
\begin{equation}\label{q3:local-spinor:reference-bounds}
 |\nabla^S_t\psi_{0,t}|+|D_t\psi_{0,t}|
       \le C\bigl((1-t)r^{-2}+r^{-1-2q_0}\bigr),
 \qquad
 \|\nabla^S_t\psi_{0,t}\|_2+\|D_t\psi_{0,t}\|_2\le C.
\end{equation}
The second estimate also uses positive local $C^2$ convergence
on the fixed cutoff annulus.  For example, the tail of its square
is bounded by a constant times
$\int_{2r_0}^{\infty}((1-t)^2r^{-2}+r^{-4q_0})\,dr$.
The reference field has finite derivative norms; it is not an
element of $\mathcal H_t$, since its weighted norm in
Equation~\eqref{q3:local-spinor:hardy} is infinite.

\begin{lemma}[A harmonic spinor with controlled correction]
\label{q3:local-spinor:harmonic-spinor}
For every $t_0\le t<1$ there exists $\zeta_t\in\mathcal H_t$ such
that
\begin{equation}\label{q3:local-spinor:projection-bound}
 \psi_t=\psi_{0,t}+\zeta_t,\qquad D_t\psi_t=0,
 \qquad
 \|\nabla^S_t\zeta_t\|_2\le\|D_t\psi_{0,t}\|_2\le C.
\end{equation}
\end{lemma}

\begin{proof}
Project $-D_t\psi_{0,t}$ orthogonally onto the closed image in
Equation~\eqref{q3:local-spinor:dirac-isometry}.  The unique preimage is
$\zeta_t$, and the projection inequality gives the norm bound.
The residual
$w=D_t(\psi_{0,t}+\zeta_t)\in L^2$
is orthogonal to $D_t\mathcal H_t$, so formal symmetry against
compact tests gives $D_tw=0$ distributionally.

We give the regularity step needed to test this equation.
In a coordinate spin trivialization,
$D_t=A^j(x)\partial_j+B(x)$, where the principal coefficients
are locally Lipschitz and $B$ is locally bounded.  The Clifford
relations make the principal symbol elliptic.  Freezing $A^j$
at a point and applying the Fourier transform gives, for a
smooth compact test vector in a sufficiently small chart,
\[
 \|v\|_{H^1}\le C\bigl(\|D_tv\|_{L^2}+\|v\|_{L^2}\bigr).
\]
Indeed the constant symbol bounds $|\xi||\widehat v(\xi)|$;
the small oscillation of $A^j$ on the chart is absorbed into
the left side, and the bounded zeroth-order coefficient
contributes the $L^2$ term.  This estimate applies to a localized
weak solution by mollification.  To see the necessary uniform
bound, the commutator of $A^j\partial_j$ with a mollifier has
kernels involving
$(A^j(x)-A^j(y))\partial_j\rho_\varepsilon(x-y)$ and
$(\partial_jA^j)(y)\rho_\varepsilon(x-y)$.
Their $L^1$ norms are uniformly bounded by the Lipschitz norm
of $A^j$.  The commutator with $B$ is bounded on $L^2$ as well.
For a compact localization $v$ of $w$, both $v$ and $D_tv$
are in $L^2$, so the mollified fields are uniformly $H^1$ on
smaller charts and their weak limit is $v$.  Hence
$w\in H^1_{\mathrm{loc}}$.  This is also the local regularity
result in \cite[Theorem~3.7]{BartnikChrusciel2005}; the present
$C^2$ metric satisfies its coefficient hypotheses.

For the fixed complete metric $q_t$, the length-space form of
Hopf--Rinow makes closed metric balls compact
\cite[Section 2.5]{BuragoBuragoIvanov2001}.  Distance from a
fixed point therefore supplies compactly supported Lipschitz
cutoffs $\chi_R$ which equal one on the ball of radius $R$,
vanish outside the ball of radius $2R$, and satisfy
$|d\chi_R|_{q_t}\le C/R$ almost everywhere.
The product $\chi_Rw$ is compactly supported $H^1$.  Since
$D_tw=0$, Equation~\eqref{q3:local-spinor:lichnerowicz} gives
\[
 \|\nabla^S_t(\chi_Rw)\|_2^2
       =\|\mathfrak c_t(d\chi_R)w\|_2^2
       \le CR^{-2}\|w\|_2^2.
\]
Letting $R\to\infty$ shows on each compact set that
$\nabla^S_tw=0$.  Its length is constant on the connected
$W$, by Kato or the metric compatibility of the connection.
The distinguished end has infinite $q_t$ volume.  Since
$w\in L^2$, this constant is zero, and $w=0$ as required.
\end{proof}

\subsection{A nonzero parallel spinor in the local limit}

The harmonic spinors have uniformly bounded corrections. We now show
that their total covariant-derivative energy tends to zero. This stronger
estimate will make the local limit parallel, while the weighted bound
on the correction will keep it nonzero.

\begin{lemma}[The finite-energy identity]
\label{q3:local-spinor:energy-lemma}
The spinors in Lemma~\ref{q3:local-spinor:harmonic-spinor} satisfy
\begin{equation}\label{q3:local-spinor:small-energy}
 \|\nabla^S_t\psi_t\|_{L^2(q_t)}^2\le C(1-t),
 \qquad t_0\le t<1.
\end{equation}
\end{lemma}

\begin{proof}
On the fields with finite derivative and Dirac norms used here,
define the continuous bilinear form
\[
 \mathcal B_t(\xi,\eta)
  =\langle\nabla^S_t\xi,\nabla^S_t\eta\rangle_{L^2}
       -\langle D_t\xi,D_t\eta\rangle_{L^2}.
\]
The polarized compact identity gives $\mathcal B_t(\xi,\eta)=0$
whenever one entry is compactly supported and both are locally
$H^1$: multiply the other entry by a compact cutoff which is
one near that support and use
Equation~\eqref{q3:local-spinor:lichnerowicz}.  Approximate
$\zeta_t$ by compact tests in $\mathcal H_t$.
Equation~\eqref{q3:local-spinor:reference-bounds} permits passage to the
limit in the mixed terms, giving
\begin{equation}\label{q3:local-spinor:mixed-identity}
 \mathcal B_t(\zeta_t,\zeta_t)=0,
 \quad \mathcal B_t(\psi_{0,t},\zeta_t)=0,
 \quad
 \mathcal B_t(\psi_t,\psi_t)
       =\mathcal B_t(\psi_{0,t},\psi_{0,t}).
\end{equation}
This use of compact approximation concerns the correction
$\zeta_t$.

To estimate the last expression, take a scalar cutoff $\beta_R$
which is one up to radius $R$ in the distinguished end and zero
beyond $2R$, with $|d\beta_R|_{q_t}\le C/R$ there; extend it
as one off the end.  The spinor $\beta_R\psi_{0,t}$ is compactly
supported.  The product rules for $\nabla^S_t$ and $D_t$ give
\begin{align*}
 0={}&\mathcal B_t(\psi_{0,t},\beta_R\psi_{0,t})\\
   ={}&\int_W\beta_R
       (|\nabla^S_t\psi_{0,t}|^2-|D_t\psi_{0,t}|^2)\,dV_{q_t}
       +\mathcal E_{t,R},
\end{align*}
where
\begin{align}
 |\mathcal E_{t,R}|
 &\le\frac C R\int_{R<r<2R}|\psi_{0,t}|
       (|\nabla^S_t\psi_{0,t}|+|D_t\psi_{0,t}|)\,dV_{q_t}
       \notag\\
 &\le C\bigl((1-t)+R^{1-2q_0}\bigr).
       \label{q3:local-spinor:annulus-error}
\end{align}
Here $|\psi_{0,t}|=1$ on this annulus for large $R$, its
volume is $O(R^3)$ uniformly in $t$, and
Equation~\eqref{q3:local-spinor:reference-bounds} supplies the two
derivative rates.  The integral on the preceding line converges
to $\mathcal B_t(\psi_{0,t},\psi_{0,t})$ by integrability of
the two squared derivatives.  Since $2q_0>1$, letting
$R\to\infty$ in Equation~\eqref{q3:local-spinor:annulus-error} yields
$|\mathcal B_t(\psi_{0,t},\psi_{0,t})|\le C(1-t)$.
Finally $D_t\psi_t=0$ and
Equation~\eqref{q3:local-spinor:mixed-identity} prove
Equation~\eqref{q3:local-spinor:small-energy}.
All explicit flux estimates involve the reference field, whose
support lies in the chosen end neighborhood; the correction was
handled by its compact approximation.
\end{proof}

\begin{proposition}[Local flatness of the limiting metric]
\label{q3:local-spinor:flatness}
The metric $q_*$ is flat on the plus interior $\mathcal P$ and
has zero curvature at each attached sphere $S_i$.
\end{proposition}

\begin{proof}
Use the fixed spin structure and bundle identifications with the
positive $C^2$ metric $q_*$ on compact subsets of $W$.
Equations~\eqref{q3:local-spinor:projection-bound} and
\eqref{q3:local-spinor:local-coercivity} give uniform local $H^1$ bounds
for $\zeta_t$.  The reference spinors $\psi_{0,t}$ converge
locally with their coefficients, so $\psi_t$ also has uniform
local $H^1$ bounds.  A diagonal subsequence as $t\uparrow1$
converges weakly in $H^1$ on every compact subdomain to a
spinor $\psi_*$; write
$\zeta_*:=\psi_*-\psi_{0,*}$.
The connection coefficients converge locally, and
Equation~\eqref{q3:local-spinor:small-energy} implies
\begin{equation}\label{q3:local-spinor:parallel-limit}
 \nabla^S_*\psi_*=0
 \quad\hbox{in distributions on }W.
\end{equation}

It remains to prove that this limit does not vanish.
Equations~\eqref{q3:local-spinor:hardy} and
\eqref{q3:local-spinor:projection-bound} give the uniform estimate
\[
 \int_{r\ge r_0}r^{-2}|\zeta_t|^2\,dV_{q_t}\le C.
\]
On every bounded end annulus, local weak lower semicontinuity
and the converging metric and bundle coefficients pass this
bound to $\zeta_*$.  Exhausting the end by such annuli gives
\begin{equation}\label{q3:local-spinor:limit-weight}
 \int_{r\ge r_0}r^{-2}|\zeta_*|^2\,dV_{q_*}\le C.
\end{equation}
If $\psi_*$ vanished identically on $\mathcal P$, then
$\zeta_*=-\psi_{0,*}$ on this end.  The reference spinor has
unit length for $r\ge2r_0$, and $q_*$ is uniformly comparable
to the Euclidean metric there, so the integral in
Equation~\eqref{q3:local-spinor:limit-weight} would be bounded below
by a positive constant times $\int_{2r_0}^{\infty}dr=\infty$.
Thus $\psi_*$ is nonzero on the plus interior.

The metric is smooth on $\mathcal P$.  The parallel equation
there first improves the local Sobolev regularity of $\psi_*$
by differentiating in smooth local trivializations, and iteration
gives a smooth parallel spinor.  Metric compatibility shows that
its length is constant; connectedness of $\mathcal P$ makes that
constant positive.  Its spin curvature therefore annihilates it.
The Clifford contraction of the spin-curvature formula, using
the first Bianchi identity, is
\[
 \sum_j\mathfrak c_*(e_j)R^S_*(e_j,Y)\psi_*
       =\tfrac12\mathfrak c_*
          \bigl(\Ric_{q_*}(Y,\cdot)^{\sharp}\bigr)\psi_*=0.
\]
Changing both curvature conventions changes the displayed sign
but leaves the conclusion unchanged.  Clifford multiplication by
a nonzero vector is invertible, since its square is minus its
squared length.  Hence $\Ric_{q_*}=0$ on $\mathcal P$.
In three dimensions the Riemann tensor is determined algebraically
by the Ricci tensor, so $q_*$ is flat there.  Its $C^2$ extension
across the joined spheres has continuous curvature, which proves
the last assertion.
\end{proof}

\subsection{The horizon area contradiction}

We finish the contradiction under the equality assumption
$|S_1|_g+|S_2|_g=16\pi m^2$.  On the plus side set
$\phi=(1+\lambda)/2$, so $q_*=\phi^4h$.
Equation~\eqref{q3:local-spinor:horizon-input} gives
$\phi=1/2$ and $\nu_h\log\phi=\kappa$ at each $S_i$.
The conformal normal is $\nu_{q_*}=\phi^{-2}\nu_h$, and the
conformal second fundamental form, in this orientation, is
\begin{equation}\label{q3:local-spinor:shape}
 \II_{q_*}
 =\phi^2\bigl(\II_h+2\nu_h(\log\phi)h|_{TS_i}\bigr)
 =\frac\kappa2h|_{TS_i}
 =8\kappa\,q_*|_{TS_i}.
\end{equation}
Both principal curvatures are consequently $8\kappa$.
By Proposition~\ref{q3:local-spinor:flatness} and the Gauss equation,
the intrinsic Gauss curvature of each sphere is $64\kappa^2$.
Although the ambient extension of $q_*$ is only $C^2$, its induced
metric on $S_i$ equals $g|_{TS_i}/16$ and is smooth: the original
metric and horizon sphere are smooth, and the attachment preserves
their tangential smooth structure. Thus the ordinary smooth
Gauss--Bonnet formula applies \cite[Section~1]{Chern1944} and gives
$64\kappa^2|S_i|_{q_*}=2\pi\chi(S_i)=4\pi$.  Since the induced metric is
$q_*|_{TS_i}=g|_{TS_i}/16$, its two-dimensional area form is
$1/16$ times the original area form.  Therefore
\begin{equation}\label{q3:local-spinor:individual-area}
 |S_i|_g=16|S_i|_{q_*}
       =16\frac{4\pi}{64\kappa^2}
       =\frac\pi{\kappa^2}=16\pi m^2,
 \qquad i=1,2.
\end{equation}
Their sum would be $32\pi m^2$, whereas the assumed equality
makes it $16\pi m^2$.  This is impossible because $m>0$.
Together with the already established non-strict inequality,
this excludes equality and proves
Theorem~\ref{q3:local-spinor:strictness}.
